% !TEX root = main.tex
% Bab Evaluasi Latensi Inferensi pada Mikrokontroler MSP430
% Ditulis dalam Bahasa Indonesia

\section{Evaluasi Latensi Inferensi pada Mikrokontroler MSP430}

Bagian ini menyajikan metodologi dan hasil evaluasi latensi inferensi model deteksi intrusi berbasis \textit{tree ensemble} pada mikrokontroler MSP430F5529. Evaluasi ini bertujuan untuk memvalidasi kelayakan penerapan model \textit{machine learning} yang telah dilatih pada perangkat keras \textit{Wireless Sensor Network} (WSN) dengan sumber daya terbatas.

\subsection{Justifikasi Pemilihan Platform MSP430}

Pemilihan mikrokontroler MSP430 sebagai target evaluasi didasarkan pada beberapa pertimbangan:

\begin{enumerate}
    \item \textbf{Relevansi dengan Arsitektur LEACH}: Dalam protokol LEACH (\textit{Low-Energy Adaptive Clustering Hierarchy}), setiap node sensor berpotensi menjadi \textit{Cluster Head} (CH) pada setiap putaran. Oleh karena itu, semua node harus memiliki kemampuan komputasi yang memadai untuk menjalankan model IDS. MSP430 merupakan platform yang umum digunakan dalam implementasi LEACH pada penelitian WSN.
    
    \item \textbf{Karakteristik Konsumsi Daya Rendah}: MSP430 dirancang khusus untuk aplikasi \textit{ultra-low power}, dengan konsumsi arus aktif 290~$\mu$A/MHz dan mode tidur hingga 1.6~$\mu$A. Karakteristik ini sesuai dengan kebutuhan node sensor yang beroperasi dengan baterai terbatas.
    
    \item \textbf{Kapasitas Memori yang Representatif}: Dengan 128~KB \textit{flash} dan 8~KB RAM, MSP430F5529 merepresentasikan kategori mikrokontroler 16-bit yang umum digunakan dalam aplikasi WSN komersial dan penelitian.
\end{enumerate}

\subsection{Metodologi Evaluasi Latensi Inferensi}

Metodologi evaluasi latensi inferensi terdiri dari empat tahap utama: kuantisasi model, generasi kode C, kompilasi untuk MSP430, dan analisis siklus instruksi.

\subsubsection{Tahap 1: Kuantisasi Model INT8}

Model \textit{tree ensemble} yang dilatih menggunakan \textit{scikit-learn} menggunakan representasi \textit{floating-point} 32-bit untuk nilai \textit{threshold} dan fitur. Untuk deployment pada mikrokontroler 16-bit tanpa FPU (\textit{Floating-Point Unit}), diperlukan kuantisasi ke format bilangan bulat 8-bit (INT8).

\paragraph{Skema Kuantisasi Min-Max}
Kuantisasi dilakukan menggunakan skema \textit{min-max scaling} per fitur:

\begin{equation}
    q = \text{round}\left(\frac{x - x_{min}}{x_{max} - x_{min}} \times 255\right)
\end{equation}

di mana $x$ adalah nilai fitur asli, $x_{min}$ dan $x_{max}$ adalah nilai minimum dan maksimum fitur pada data kalibrasi, dan $q$ adalah nilai terkuantisasi dalam rentang [0, 255].

Untuk setiap fitur $i$, disimpan dua parameter kuantisasi:
\begin{itemize}
    \item \textbf{Faktor Skala} ($s_i$): Disimpan sebagai bilangan bulat 16-bit dalam format \textit{fixed-point} Q8.8
    \item \textbf{Titik Nol} ($z_i$): Disimpan sebagai bilangan bulat 8-bit \textit{unsigned}
\end{itemize}

Proses kuantisasi pada saat inferensi dilakukan dengan operasi:
\begin{equation}
    q_i = \text{clamp}\left(\left\lfloor\frac{x_i \times s_i}{256}\right\rfloor + z_i, 0, 255\right)
\end{equation}

\paragraph{Kuantisasi Threshold Decision Tree}
Setiap \textit{threshold} pada node keputusan juga dikuantisasi menggunakan parameter yang sama dengan fitur terkait:

\begin{equation}
    t_q = \text{round}(t \times s_{f} / 256 + z_{f})
\end{equation}

di mana $t$ adalah \textit{threshold} asli, $f$ adalah indeks fitur yang digunakan pada node tersebut, dan $t_q$ adalah \textit{threshold} terkuantisasi.

\subsubsection{Tahap 2: Pembatasan Kedalaman Pohon}

Untuk memenuhi batasan memori MSP430 (128~KB \textit{flash}), kedalaman pohon dibatasi pada 6 level. Pembatasan ini memberikan keseimbangan antara kompleksitas model dan penggunaan memori:

\begin{itemize}
    \item Jumlah node maksimum per pohon: $2^6 - 1 = 63$ node
    \item Ukuran per node: 6 byte (indeks fitur, \textit{threshold}, \textit{pointer} anak kiri dan kanan)
    \item Memori per pohon: maksimal $63 \times 6 = 378$ byte
\end{itemize}

Ekstraksi pohon dilakukan menggunakan \textit{Breadth-First Search} (BFS) dengan penghentian paksa pada kedalaman maksimum. Node yang dipotong pada kedalaman maksimum dikonversi menjadi node daun dengan kelas mayoritas.

\paragraph{Format Node Terkuantisasi}
Setiap node pohon keputusan direpresentasikan dalam format 6 byte:

\begin{table}[htbp]
\centering
\caption{Format Node Pohon Keputusan Terkuantisasi}
\label{tab:node_format}
\begin{tabular}{clc}
\hline
\textbf{Offset} & \textbf{Field} & \textbf{Ukuran} \\
\hline
0 & Indeks fitur (255 = daun) & 1 byte \\
1 & Threshold / ID kelas & 1 byte \\
2-3 & Pointer anak kiri & 2 byte \\
4-5 & Pointer anak kanan & 2 byte \\
\hline
\end{tabular}
\end{table}

\subsubsection{Tahap 3: Generasi Kode C untuk MSP430}

Proses generasi kode C menghasilkan file-file berikut untuk setiap model:

\begin{enumerate}
    \item \texttt{model\_config.h}: Definisi konfigurasi model (jumlah pohon, fitur, kelas)
    \item \texttt{quantization.h}: Parameter kuantisasi (faktor skala dan titik nol)
    \item \texttt{tree\_data.h}: Data pohon terkuantisasi dalam format array konstan
    \item \texttt{inference.c/h}: Implementasi fungsi inferensi
    \item \texttt{main.c}: Program \textit{benchmark} untuk pengukuran waktu
    \item \texttt{Makefile}: Konfigurasi kompilasi untuk MSP430
\end{enumerate}

\paragraph{Implementasi Fungsi Inferensi}
Fungsi inferensi terdiri dari tiga komponen utama:

\begin{enumerate}
    \item \textbf{Kuantisasi Fitur} (\texttt{quantize\_features}): Mengkonversi 16 fitur \textit{floating-point} ke format UINT8.
    
    \item \textbf{Traversal Pohon} (\texttt{tree\_predict}): Menelusuri pohon dari akar ke daun berdasarkan perbandingan fitur dengan \textit{threshold}.
    
    \item \textbf{Voting Ensemble} (\texttt{ensemble\_predict}): Mengumpulkan prediksi dari semua pohon dan menentukan kelas mayoritas.
\end{enumerate}

Pseudocode untuk traversal pohon ditunjukkan pada Algoritma~\ref{alg:tree_predict}.

\begin{algorithm}[htbp]
\caption{Traversal Pohon Keputusan Terkuantisasi}
\label{alg:tree_predict}
\begin{algorithmic}[1]
\REQUIRE tree\_data: array data node, features: array fitur terkuantisasi
\ENSURE predicted\_class: kelas hasil prediksi
\STATE node\_idx $\leftarrow$ 0
\WHILE{true}
    \STATE feature\_idx $\leftarrow$ tree\_data[node\_idx $\times$ 6]
    \STATE threshold $\leftarrow$ tree\_data[node\_idx $\times$ 6 + 1]
    \IF{feature\_idx = 255}
        \RETURN threshold \COMMENT{Node daun, threshold berisi ID kelas}
    \ENDIF
    \IF{features[feature\_idx] $\leq$ threshold}
        \STATE node\_idx $\leftarrow$ tree\_data[node\_idx $\times$ 6 + 2:3] \COMMENT{Anak kiri}
    \ELSE
        \STATE node\_idx $\leftarrow$ tree\_data[node\_idx $\times$ 6 + 4:5] \COMMENT{Anak kanan}
    \ENDIF
\ENDWHILE
\end{algorithmic}
\end{algorithm}

\subsubsection{Tahap 4: Kompilasi dan Analisis Siklus}

\paragraph{Kompilasi untuk MSP430}
Kode C dikompilasi menggunakan \texttt{msp430-elf-gcc} dengan opsi optimisasi \texttt{-O2}:

\begin{verbatim}
msp430-elf-gcc -mmcu=msp430f5529 -O2 -Wall -c -o inference.o inference.c
msp430-elf-gcc -mmcu=msp430f5529 -o inference.elf main.o inference.o
\end{verbatim}

\paragraph{Metodologi Penghitungan Siklus}
Latensi inferensi dihitung berdasarkan analisis siklus instruksi dari kode hasil kompilasi. Penghitungan menggunakan informasi waktu eksekusi instruksi dari \textit{MSP430x5xx Family User's Guide} (SLAU208Q).

Tabel~\ref{tab:msp430_cycles} menunjukkan jumlah siklus untuk kategori instruksi utama pada MSP430:

\begin{table}[htbp]
\centering
\caption{Jumlah Siklus Instruksi MSP430}
\label{tab:msp430_cycles}
\begin{tabular}{lcc}
\hline
\textbf{Kategori Instruksi} & \textbf{Contoh} & \textbf{Siklus} \\
\hline
Operasi register-register & MOV, ADD, CMP & 1 \\
Operasi dengan memori & MOV @Rn, Rd & 2--3 \\
Lompatan kondisional & JNE, JZ, JC & 2 \\
Panggilan fungsi & CALL & 4 \\
Kembali dari fungsi & RET & 3 \\
Push/Pop stack & PUSH, POP & 2--3 \\
\hline
\end{tabular}
\end{table}

\paragraph{Model Estimasi Siklus}
Total siklus inferensi dihitung berdasarkan tiga komponen:

\begin{equation}
    C_{total} = C_{quant} + C_{trees} + C_{voting}
\end{equation}

di mana:
\begin{itemize}
    \item $C_{quant}$: Siklus untuk kuantisasi fitur
    \item $C_{trees}$: Siklus untuk traversal semua pohon
    \item $C_{voting}$: Siklus untuk \textit{voting} dan penentuan kelas
\end{itemize}

\textbf{Kuantisasi Fitur:}
Untuk setiap fitur, operasi kuantisasi memerlukan:
\begin{itemize}
    \item Load faktor skala dari memori: 3 siklus
    \item Perkalian: 8 siklus (emulasi software)
    \item Shift dan penjumlahan: 4 siklus
    \item Clamping: 6 siklus
    \item Store hasil: 4 siklus
\end{itemize}

Total per fitur: $\approx$35 siklus. Untuk 16--18 fitur: $C_{quant} \approx 650$ siklus.

\textbf{Traversal Pohon:}
Untuk setiap node yang dikunjungi:
\begin{itemize}
    \item Load indeks fitur: 3 siklus
    \item Load threshold: 3 siklus
    \item Perbandingan: 2 siklus
    \item Load pointer anak: 4 siklus
    \item Lompatan kondisional: 2 siklus
    \item Overhead loop: 3 siklus
\end{itemize}

Total per node: $\approx$25 siklus. Dengan kedalaman rata-rata 4 node (setengah dari kedalaman maksimum 6), dan overhead panggilan fungsi 10 siklus:

\begin{equation}
    C_{per\_tree} = 4 \times 25 + 10 = 110 \text{ siklus}
\end{equation}

\textbf{Voting Ensemble:}
\begin{itemize}
    \item Inisialisasi array vote: $n_{classes} \times 2$ siklus
    \item Increment vote per pohon: $n_{trees} \times 8$ siklus
    \item Pencarian maksimum: $n_{classes} \times 5$ siklus
\end{itemize}

\subsection{Spesifikasi Platform Target}

\begin{table}[htbp]
\centering
\caption{Spesifikasi Mikrokontroler MSP430F5529}
\label{tab:msp430_specs_id}
\begin{tabular}{ll}
\hline
\textbf{Parameter} & \textbf{Nilai} \\
\hline
Arsitektur & 16-bit RISC \\
Frekuensi Clock & 25 MHz \\
Memori Flash & 128 KB \\
Memori SRAM & 8 KB \\
Arus Aktif & 290 $\mu$A/MHz \\
Mode Daya Rendah (LPM3) & 1.6 $\mu$A \\
Tegangan Operasi & 1.8V -- 3.6V \\
Peripheral & Timer, UART, SPI, I2C, ADC \\
Aplikasi Tipikal & Node sensor WSN, LEACH CH \\
\hline
\end{tabular}
\end{table}

\subsection{Hasil Evaluasi Latensi Inferensi}

Evaluasi dilakukan pada empat model \textit{tree ensemble} yang dilatih tanpa \textit{feature engineering} dan tanpa \textit{oversampling}:

\begin{enumerate}
    \item \textbf{Extra Trees}: 5 pohon dengan kedalaman maksimum 6
    \item \textbf{Random Forest}: 5 pohon dengan kedalaman maksimum 6
    \item \textbf{Gradient Boosting}: 20 pohon dengan kedalaman maksimum 6
    \item \textbf{HistGradient Boosting}: 20 pohon dengan kedalaman maksimum 6
\end{enumerate}

\subsubsection{Penggunaan Memori}

\begin{table}[htbp]
\centering
\caption{Penggunaan Memori pada MSP430F5529}
\label{tab:msp430_memory_id}
\begin{tabular}{lcccc}
\hline
\textbf{Model} & \textbf{Pohon} & \textbf{Flash (byte)} & \textbf{Flash (\%)} & \textbf{RAM (\%)} \\
\hline
Extra Trees & 5 & 10.893 & 8,31\% & 3,76\% \\
Random Forest & 5 & 10.895 & 8,31\% & 3,76\% \\
Gradient Boosting & 20 & 20.247 & 15,45\% & 3,76\% \\
HistGradient Boosting & 20 & 8.335 & 6,36\% & 3,76\% \\
\hline
\end{tabular}
\end{table}

Hasil menunjukkan bahwa semua model dapat dimuat dalam memori flash MSP430F5529 dengan penggunaan kurang dari 20\%. Hal ini menyisakan ruang yang cukup untuk \textit{protocol stack} LEACH, driver radio, dan kode aplikasi lainnya.

\subsubsection{Rincian Jumlah Siklus}

\begin{table}[htbp]
\centering
\caption{Rincian Jumlah Siklus per Komponen}
\label{tab:msp430_cycles_detail}
\begin{tabular}{lcccc}
\hline
\textbf{Komponen} & \textbf{Extra Trees} & \textbf{Random Forest} & \textbf{Gradient} & \textbf{HistGradient} \\
\hline
Kuantisasi Fitur & 650 & 650 & 650 & 650 \\
Traversal Pohon & 550 & 550 & 2.200 & 2.200 \\
Overhead Ensemble & 70 & 70 & 190 & 190 \\
\textbf{Total Siklus} & \textbf{1.270} & \textbf{1.270} & \textbf{3.040} & \textbf{3.040} \\
\hline
\end{tabular}
\end{table}

Dari tabel di atas terlihat bahwa:
\begin{itemize}
    \item Komponen kuantisasi fitur konstan untuk semua model karena jumlah fitur sama (18 fitur).
    \item Komponen traversal pohon berbanding lurus dengan jumlah pohon: model dengan 20 pohon memerlukan 4$\times$ siklus dibanding model dengan 5 pohon.
    \item Overhead ensemble meningkat linear dengan jumlah pohon karena operasi lookup pointer dan increment vote.
\end{itemize}

\subsubsection{Hasil Waktu Inferensi}

Waktu inferensi dihitung dari jumlah siklus dibagi frekuensi clock 25 MHz:

\begin{equation}
    t_{inference} = \frac{C_{total}}{f_{clock}} = \frac{C_{total}}{25 \times 10^6} \text{ detik}
\end{equation}

\begin{table}[htbp]
\centering
\caption{Hasil Latensi Inferensi pada MSP430F5529 @ 25 MHz}
\label{tab:msp430_latency_id}
\begin{tabular}{lccc}
\hline
\textbf{Model} & \textbf{Latensi ($\mu$s)} & \textbf{Latensi (ms)} & \textbf{Throughput (inf/s)} \\
\hline
Extra Trees & 50,8 & 0,0508 & 19.685 \\
Random Forest & 50,8 & 0,0508 & 19.685 \\
Gradient Boosting & 121,6 & 0,1216 & 8.224 \\
HistGradient Boosting & 121,6 & 0,1216 & 8.224 \\
\hline
\end{tabular}
\end{table}

\subsubsection{Estimasi Konsumsi Energi}

Konsumsi energi per inferensi dihitung berdasarkan karakteristik daya MSP430F5529:

\begin{equation}
    E_{inference} = I_{active} \times V_{supply} \times t_{inference}
\end{equation}

Dengan arus aktif 290~$\mu$A/MHz pada 25~MHz (total 7,25~mA) dan tegangan 3,3V:

\begin{table}[htbp]
\centering
\caption{Konsumsi Energi per Inferensi}
\label{tab:msp430_energy_id}
\begin{tabular}{lcc}
\hline
\textbf{Model} & \textbf{Waktu ($\mu$s)} & \textbf{Energi ($\mu$J)} \\
\hline
Extra Trees & 50,8 & 1,22 \\
Random Forest & 50,8 & 1,22 \\
Gradient Boosting & 121,6 & 2,91 \\
HistGradient Boosting & 121,6 & 2,91 \\
\hline
\end{tabular}
\end{table}

Sebagai perbandingan, transmisi radio menggunakan transceiver CC2420 (umum digunakan dalam WSN) mengkonsumsi sekitar 1~mJ per paket~\cite{polastre2005versatile}. Dengan demikian, energi inferensi hanya mewakili sekitar 0,1--0,3\% dari energi transmisi radio.

\subsection{Analisis dan Pembahasan}

\subsubsection{Kesesuaian dengan Batasan Waktu Protokol LEACH}

Dalam protokol LEACH, setiap putaran terdiri dari:
\begin{itemize}
    \item \textbf{Fase Setup} ($\approx$1--2 detik): Pemilihan CH, pembentukan cluster, penjadwalan TDMA
    \item \textbf{Fase Steady-state} ($\approx$18--19 detik): Pengiriman data dari anggota cluster ke CH
\end{itemize}

Selama fase \textit{steady-state}, CH menerima data dari 10--20 anggota cluster dengan interval 1--2 detik per anggota. Dalam skenario terburuk dengan 10 paket/detik:

\begin{equation}
    t_{packet\_interval} = 100 \text{ ms}
\end{equation}

Rasio margin waktu:
\begin{equation}
    \text{Margin} = \frac{t_{packet\_interval}}{t_{inference}} = \frac{100 \text{ ms}}{0,05 \text{ ms}} \approx 2000\times
\end{equation}

Margin waktu sebesar 2000$\times$ mengindikasikan bahwa sistem IDS dapat beroperasi secara \textit{real-time} tanpa mengganggu operasi normal protokol LEACH.

\subsubsection{Perbandingan antar Model}

\begin{enumerate}
    \item \textbf{Extra Trees dan Random Forest}: Kedua model menunjukkan performa identik dalam hal latensi (50,8~$\mu$s) dan penggunaan memori ($\approx$11~KB). Keduanya menggunakan 5 pohon dengan struktur serupa. Model-model ini direkomendasikan untuk aplikasi yang mengutamakan latensi minimum.
    
    \item \textbf{Gradient Boosting}: Memerlukan waktu inferensi 2,4$\times$ lebih lama (121,6~$\mu$s) karena menggunakan 20 pohon. Penggunaan memori flash tertinggi (20~KB) disebabkan oleh kompleksitas pohon yang lebih besar. Namun, latensi tetap jauh di bawah batasan \textit{real-time}.
    
    \item \textbf{HistGradient Boosting}: Meskipun menggunakan 20 pohon seperti Gradient Boosting, penggunaan memori flash jauh lebih kecil (8~KB). Hal ini karena implementasi menggunakan struktur pohon yang lebih sederhana. Model ini cocok untuk skenario dengan batasan memori ketat.
\end{enumerate}

\subsubsection{Implikasi untuk Deployment WSN}

Hasil evaluasi mengkonfirmasi kelayakan deployment model IDS berbasis \textit{tree ensemble} pada node sensor WSN:

\begin{enumerate}
    \item \textbf{Efisiensi Memori}: Semua model menggunakan kurang dari 20\% kapasitas flash, menyisakan ruang yang cukup untuk:
    \begin{itemize}
        \item Stack protokol LEACH ($\approx$10--15~KB)
        \item Driver radio CC2420 ($\approx$5~KB)
        \item Kode aplikasi sensor ($\approx$5--10~KB)
    \end{itemize}
    
    \item \textbf{Kapabilitas Real-time}: Latensi sub-milidetik memungkinkan inspeksi per-paket tanpa buffering, memberikan deteksi ancaman segera.
    
    \item \textbf{Efisiensi Energi}: Konsumsi energi inferensi (1--3~$\mu$J) sangat kecil dibanding operasi radio, sehingga tidak berdampak signifikan pada masa pakai baterai.
    
    \item \textbf{Oversampling Tidak Diperlukan}: Mengingat model tanpa \textit{oversampling} telah mencapai akurasi 99\%+ pada dataset WSN-DS, kompleksitas tambahan dari teknik \textit{oversampling} tidak memberikan manfaat deployment.
\end{enumerate}

\subsubsection{Keterbatasan Evaluasi}

Beberapa keterbatasan dari metodologi evaluasi ini perlu diperhatikan:

\begin{enumerate}
    \item \textbf{Estimasi Siklus Statis}: Penghitungan siklus dilakukan secara statis berdasarkan analisis instruksi, bukan pengukuran langsung pada perangkat keras. Variasi akibat \textit{cache miss} atau \textit{interrupt} tidak diperhitungkan.
    
    \item \textbf{Pembatasan Kedalaman Pohon}: Pembatasan kedalaman maksimum 6 level dapat mempengaruhi akurasi model. Evaluasi dampak akurasi dari pembatasan ini tidak termasuk dalam cakupan penelitian ini.
    
    \item \textbf{Kuantisasi INT8}: Konversi dari \textit{floating-point} ke INT8 dapat menyebabkan degradasi akurasi. Analisis kuantitatif degradasi akurasi memerlukan evaluasi terpisah.
\end{enumerate}

\subsection{Kesimpulan Evaluasi MSP430}

Evaluasi latensi inferensi pada mikrokontroler MSP430F5529 menunjukkan bahwa:

\begin{enumerate}
    \item Model \textit{tree ensemble} terkuantisasi INT8 dapat dieksekusi dengan latensi 50,8--121,6~$\mu$s pada frekuensi 25~MHz.
    
    \item Semua model yang dievaluasi memenuhi batasan memori flash 128~KB dengan penggunaan 6--16\%.
    
    \item Throughput inferensi (8.000--20.000 inferensi/detik) melebihi kebutuhan tipikal WSN sebesar 100--1000$\times$.
    
    \item Konsumsi energi per inferensi (1--3~$\mu$J) dapat diabaikan dibanding energi transmisi radio.
    
    \item Deployment IDS berbasis \textit{tree ensemble} pada Cluster Head LEACH layak dilakukan tanpa mengorbankan performa protokol atau masa pakai baterai.
\end{enumerate}

Temuan ini mendukung kesimpulan bahwa teknik \textit{oversampling} tidak diperlukan untuk mencapai performa deteksi intrusi yang baik pada platform WSN dengan sumber daya terbatas. Model yang dilatih tanpa \textit{oversampling} memberikan akurasi tinggi dengan kompleksitas komputasi yang efisien.
